% Numbers come from report/generated/annotation_macros.tex (src/analyze_annotations.py).
\paragraph{Agreement between people.} All three annotators completed all 60 items (one Q2 answer was left blank). Inter-annotator agreement on trust is Krippendorff's $\alpha$\,=\,\AlphaTrust{} (ordinal), and pairwise quadratic-weighted $\kappa$ ranges over \PairQWKRange{}, i.e.\ moderate to substantial agreement. Agreement on whether an explanation ``makes sense'' is lower ($\alpha$\,=\,\AlphaSense{}, nominal), so judging the medical plausibility of a reason is harder for non-experts than giving an overall trust rating. Annotator C was the most sceptical, and the mean trust over all ratings is \MeanTrustHuman{} on the 1--5 scale.

\paragraph{Do explanations help people spot wrong predictions?} Correct predictions received a mean human trust of \TrustCorrect{}, and incorrect ones \TrustIncorrect{} (Mann--Whitney $p$\,=\,\PTrustMW{}). Averaged trust separates correct from incorrect predictions with ROC AUC \AUCHuman{} (95\% bootstrap CI \AUCLo{}--\AUCHi{}; 0.5 is chance), and explanations were judged to make sense for \SenseCorrect{} of correct but only \SenseIncorrect{} of incorrect predictions. On the face of it this supports the paper's claim. However, human trust correlates almost perfectly with the model confidence printed next to each prediction (Spearman $\rho$\,=\,\RhoHumanConf{}), and confidence \emph{alone} separates correct from incorrect predictions just as well (AUC \AUCConf{}). Once confidence is regressed out, the remaining variation in human trust carries no signal about correctness (AUC \AUCResid{}, $p$\,=\,\PResid{}), and trust correlates only weakly with LIME's local fidelity ($\rho$\,=\,\RhoHumanRtwo{}). Our raters could tell good predictions from bad, but we find no evidence that LIME's explanations contributed beyond the confidence score. Because every item showed both, the design cannot rule out a smaller effect; an explanation-versus-no-explanation condition would be needed to measure it (\S\ref{sec:concl}).

\paragraph{LLM-as-judge.} The judge (\LLMModel{}, three repeats) is almost perfectly self-consistent ($\alpha$\,=\,\LLMSelfAlpha{}), and it ranks items much as people do: Spearman $\rho$\,=\,\LLMvsHumanRho{} with the mean human rating, and quadratic-weighted $\kappa$\,=\,\LLMvsHumanQWK{} with the human median. That is within the range of agreement between the humans themselves (\PairQWKRange{}). As a rater, however, it is systematically lenient. Its mean trust is \MeanTrustLLM{} against \MeanTrustHuman{} for humans, it never gave a rating below 3 (humans gave 1--2 in 18\% of ratings), and it answered ``Yes, the explanation makes sense'' for all 180 judgements, against 63\% for humans. Agreement with the human Q2 majority is therefore $\kappa$\,=\,\LLMvsHumanSenseKappa{} by construction, and adding the judge as a fourth rater \emph{lowers} the group's $\alpha$ from \AlphaTrust{} to \AlphaTrustWithLLM{}. Its trust separates correct from incorrect predictions less well than the humans' (AUC \AUCLLM{} vs.\ \AUCHuman{}; mean \LLMTrustCorrect{} vs.\ \LLMTrustIncorrect{}). Like the humans, it mostly follows the displayed confidence ($\rho$\,=\,\RhoLLMConf{}), and beyond confidence it carries no signal (AUC \AUCLLMResid{}). Its stated reasons are generic (e.g.\ ``the explanation aligns with the patient's clinical attributes that are known risk factors''). This confirms the risk anticipated in our proposal: an LLM judge accepts fluent, plausible-sounding explanations rather than checking them. It is a cheap, reproducible proxy for the \emph{ranking} of human trust, but not a substitute for human judgement of whether an explanation is sound.
